%
\begin{isabellebody}%
\setisabellecontext{Week{\isacharunderscore}{\kern0pt}{\isadigit{2}}}%
%
\isadelimtheory
%
\endisadelimtheory
%
\isatagtheory
%
\endisatagtheory
{\isafoldtheory}%
%
\isadelimtheory
%
\endisadelimtheory
%
\begin{isamarkuptext}%
\ExerciseSheet{7}{}%
\end{isamarkuptext}\isamarkuptrue%
%
\begin{isamarkuptext}%
\Exercise{Fold function}
  The fold function is a very generic function, that can be used to express 
  multiple other interesting functions over lists.%
\end{isamarkuptext}\isamarkuptrue%
%
\begin{isamarkuptext}%
Have a look at Isabelle/HOL's standard function \isa{fold}.%
\end{isamarkuptext}\isamarkuptrue%
\isacommand{thm}\isamarkupfalse%
\ fold{\isachardot}{\kern0pt}simps%
\begin{isamarkuptext}%
Write a function to compute the sum of the elements of a list. 
  Define two versions, one direct recursive definition, and one using fold.
  Show that both are equal.

  Hint: use automation!%
\end{isamarkuptext}\isamarkuptrue%
\isacommand{fun}\isamarkupfalse%
\ list{\isacharunderscore}{\kern0pt}sum\ {\isacharcolon}{\kern0pt}{\isacharcolon}{\kern0pt}\ {\isachardoublequoteopen}nat\ list\ {\isasymRightarrow}\ nat{\isachardoublequoteclose}\ \isanewline
\ \ \isanewline
\isacommand{definition}\isamarkupfalse%
\ list{\isacharunderscore}{\kern0pt}sum{\isacharprime}{\kern0pt}\ {\isacharcolon}{\kern0pt}{\isacharcolon}{\kern0pt}\ {\isachardoublequoteopen}nat\ list\ {\isasymRightarrow}\ nat{\isachardoublequoteclose}\isanewline
\ \ %
\isadelimproof
%
\endisadelimproof
%
\isatagproof
%
\endisatagproof
{\isafoldproof}%
%
\isadelimproof
\isanewline
%
\endisadelimproof
\isacommand{lemma}\isamarkupfalse%
\ {\isachardoublequoteopen}list{\isacharunderscore}{\kern0pt}sum\ xs\ {\isacharequal}{\kern0pt}\ list{\isacharunderscore}{\kern0pt}sum{\isacharprime}{\kern0pt}\ xs{\isachardoublequoteclose}\isanewline
\ \ %
\isadelimproof
%
\endisadelimproof
%
\isatagproof
%
\endisatagproof
{\isafoldproof}%
%
\isadelimproof
%
\endisadelimproof
%
\begin{isamarkuptext}%
\Exercise{Tail recursive \isa{reverse}}
In Isabelle, there is the reverse function \isa{rev}, which, given a list, computes another list
with the same elements but in reverse order.
Take a loot into its definition:%
\end{isamarkuptext}\isamarkuptrue%
\isacommand{thm}\isamarkupfalse%
\ rev{\isachardot}{\kern0pt}simps%
\begin{isamarkuptext}%
Recall tail recursion: a function is tail recursive if in all recursive calls it does not call 
  anything after itself\footnote{See \url{https://en.wikipedia.org/wiki/Tail_call.} for a quick review.}.
  For instance, the implementaion of \isa{rev} in Isabelle is not tail recursive, because it calls
  \isa{{\isacharparenleft}{\kern0pt}{\isacharat}{\kern0pt}{\isacharparenright}{\kern0pt}} after it calls \isa{rev} in the second equation.
  
  Write an alternative implementation of the reverse function that is tail recursive.

  Hint: recall that you can most of the times derive a tail recursive function using an accumulator
        argument.%
\end{isamarkuptext}\isamarkuptrue%
\isanewline
\isanewline
\isanewline
\isacommand{fun}\isamarkupfalse%
\ rev{\isacharunderscore}{\kern0pt}tr{\isacharcolon}{\kern0pt}{\isacharcolon}{\kern0pt}{\isachardoublequoteopen}{\isacharprime}{\kern0pt}a\ list\ {\isasymRightarrow}\ {\isacharprime}{\kern0pt}a\ list{\isachardoublequoteclose}\ \isakeyword{where}%
\begin{isamarkuptext}%
Show that the two implementations are equivalent.%
\end{isamarkuptext}\isamarkuptrue%
%
\isadelimproof
%
\endisadelimproof
%
\isatagproof
%
\endisatagproof
{\isafoldproof}%
%
\isadelimproof
\ \isanewline
%
\endisadelimproof
\isacommand{lemma}\isamarkupfalse%
\ {\isachardoublequoteopen}rev{\isacharunderscore}{\kern0pt}tr\ xs\ {\isacharequal}{\kern0pt}\ rev\ xs{\isachardoublequoteclose}%
\isadelimproof
%
\endisadelimproof
%
\isatagproof
%
\endisatagproof
{\isafoldproof}%
%
\isadelimproof
%
\endisadelimproof
%
\isadelimtheory
%
\endisadelimtheory
%
\isatagtheory
%
\endisatagtheory
{\isafoldtheory}%
%
\isadelimtheory
%
\endisadelimtheory
%
\end{isabellebody}%
\endinput
%:%file=Week_2.tex%:%
%:%19=6%:%
%:%23=8%:%
%:%24=9%:%
%:%25=10%:%
%:%29=13%:%
%:%31=14%:%
%:%32=14%:%
%:%34=17%:%
%:%35=18%:%
%:%36=19%:%
%:%37=20%:%
%:%38=21%:%
%:%40=26%:%
%:%41=26%:%
%:%42=27%:%
%:%42=32%:%
%:%43=33%:%
%:%44=33%:%
%:%45=34%:%
%:%56=42%:%
%:%59=43%:%
%:%60=43%:%
%:%61=44%:%
%:%76=48%:%
%:%77=49%:%
%:%78=50%:%
%:%79=51%:%
%:%81=54%:%
%:%82=54%:%
%:%84=57%:%
%:%85=58%:%
%:%86=59%:%
%:%87=60%:%
%:%88=61%:%
%:%89=62%:%
%:%90=63%:%
%:%91=64%:%
%:%92=65%:%
%:%94=74%:%
%:%95=75%:%
%:%96=76%:%
%:%97=77%:%
%:%98=77%:%
%:%100=82%:%
%:%113=105%:%
%:%116=106%:%
%:%117=106%:%
